\section{Complete proofs of the Guinand--Weil interfaces}
\label{sec:complete-guinand-weil-proofs}

This section expands morphism records \texttt{0012--0017}.  It keeps
Guinand's RH hypothesis and half-line cosine transform separate from Weil's
general Hecke-character formula until an explicit common subspace and every
constant have been written down.  The primary inputs are Guinand's
Theorems~3 and~$3'$ and Weil's conditions (A), (B), and formula~(11)
\cite[pp.~112--115]{guinand1948}
\cite[Collected Papers II, pp.~52--58]{weil1952}.

\subsection{Weil's multiplicative test coordinate}

\begin{theorem}[Exact logarithmic Mellin test-function bijection]
\label{thm:weil-log-mellin-bijection-complete}
Let $F:\R\to\C$ satisfy Weil's condition (A), and for some $b>0$ let
\begin{equation}
 F(x),F'(x)=O\!\left(e^{-(1/2+b)|x|}\right)
 \qquad(|x|\to\infty)
 \label{eq:weil-B-restated}
\end{equation}
away from the finitely many derivative-jump points, as in Weil's condition
(B).  Define
\begin{equation}
 \mathcal B(F)(v)=v^{-1/2}F(\log v),\qquad v>0.
 \label{eq:B-map-complete}
\end{equation}
Then $\mathcal B$ is injective and has the explicit inverse
\begin{equation}
 \mathcal B^{-1}(\phi)(x)=e^{x/2}\phi(e^x).
 \label{eq:B-inverse-complete}
\end{equation}
Consequently it is a bijection from Weil's additive test space onto its
explicitly transported multiplicative test space, and every fibre is a
singleton.  At points away from transported jumps,
\begin{equation}
 \phi'(v)=v^{-3/2}
 \left(F'(\log v)-\frac12F(\log v)\right).
 \label{eq:B-derivative-complete}
\end{equation}
If $a_j$ is a jump location of $F$ or $F'$, its location becomes $e^{a_j}$;
the one-sided values and their averaging convention are transported without
merging.  The endpoint estimates are
\begin{align}
 \phi(v)&=O(v^b),&\phi'(v)&=O(v^{b-1})&&(v\downarrow0),\label{eq:B-zero-decay}\\
 \phi(v)&=O(v^{-1-b}),&\phi'(v)&=O(v^{-2-b})&&(v\to\infty).
 \label{eq:B-infinity-decay}
\end{align}
Finally, wherever Weil's transform is defined,
\begin{equation}
 \int_0^\infty \phi(v)v^{s-1}\,dv
 =\int_\R F(x)e^{(s-1/2)x}\,dx=\Phi(s).
 \label{eq:B-mellin-complete}
\end{equation}
\end{theorem}

\begin{proof}
Substituting \eqref{eq:B-map-complete} into
\eqref{eq:B-inverse-complete} gives
$e^{x/2}e^{-x/2}F(x)=F(x)$; substituting
\eqref{eq:B-inverse-complete} into \eqref{eq:B-map-complete} gives
$v^{-1/2}v^{1/2}\phi(v)=\phi(v)$.  Thus the maps are two-sided inverses.
Differentiating $v^{-1/2}F(\log v)$ gives
\eqref{eq:B-derivative-complete}.  Since $x\mapsto e^x$ is strictly
increasing and bijective, every finite jump location and both of its one-sided
germs move bijectively from $a_j$ to $e^{a_j}$.

For $v\to\infty$, put $x=\log v>0$.  Equation
\eqref{eq:weil-B-restated} gives
$F(\log v),F'(\log v)=O(v^{-1/2-b})$, which in
\eqref{eq:B-map-complete} and \eqref{eq:B-derivative-complete} gives
\eqref{eq:B-infinity-decay}.  For $v\downarrow0$, $x=\log v<0$ and
$e^{-(1/2+b)|x|}=e^{(1/2+b)x}=v^{1/2+b}$, giving
\eqref{eq:B-zero-decay}.  Finally, $v=e^x$ and $dv=e^x dx$ give
\[
 \phi(v)v^{s-1}dv
 =e^{-x/2}F(x)e^{(s-1)x}e^x dx
 =F(x)e^{(s-1/2)x}dx.
\]
This proves \eqref{eq:B-mellin-complete}.  The transform and conditions are
Weil's literal ones
\cite[Collected Papers II, pp.~52 and~55--58]{weil1952}; no endpoint decay,
jump, or averaging datum has been normalized away.
\end{proof}

\subsection{The exact Guinand--Weil intersection}

For $1<p\leq2$, let $\mathcal G_p$ denote the half-line functions satisfying
Guinand's Theorem~3 hypotheses: $f$ is an integral in Guinand's terminology,
$f(x)\to0$, $xf'(x)\in L^p(0,\infty)$, and its cosine transform is
\begin{equation}
 g(y)=\sqrt{\frac2\pi}\int_0^\infty f(x)\cos(xy)\,dx.
 \label{eq:guinand-transform-B}
\end{equation}
Let $\mathcal G_p^W$ be the subset for which the scaled even extension below
satisfies Weil's (A) and (B), and let $\mathcal W_{\mathrm{even}}^{G,p}$ be the
even Weil tests whose positive-half restriction, multiplied by
$\sqrt{2\pi}$, lies in $\mathcal G_p$.  These definitions retain the actual
intersection; they do not assert equality of the two full test spaces.

\begin{theorem}[Scaled even-extension bijection and boundary removal]
\label{thm:guinand-weil-even-complete}
The map
\begin{equation}
 \mathcal E:\mathcal G_p^W\longrightarrow
 \mathcal W_{\mathrm{even}}^{G,p},\qquad
 \mathcal E(f)(x)=\frac{f(|x|)}{\sqrt{2\pi}},
 \label{eq:even-extension-complete}
\end{equation}
is a bijection with inverse
$f(t)=\sqrt{2\pi}F(t)$ for $t\geq0$.  Its fibres are singletons.  If
$F=\mathcal E(f)$, then
\begin{equation}
 \Phi_F(\tfrac12+iy)
 =\int_\R F(x)e^{iyx}\,dx
 =\sqrt{\frac2\pi}\int_0^\infty f(x)\cos(xy)\,dx
 =g(y).
 \label{eq:even-transform-complete}
\end{equation}
Moreover, Weil's condition (B) implies
$T^{1/2}f(T)\to0$, so the boundary term in Guinand's Theorem~3 is removable
on this intersection by his Theorem~$3'$.
\end{theorem}

\begin{proof}
The image in \eqref{eq:even-extension-complete} is even.  Restriction to
$[0,\infty)$ followed by multiplication by $\sqrt{2\pi}$ is visibly a
two-sided inverse, proving the bijection and singleton fibres.  Evenness gives
\begin{align*}
 \int_\R F(x)e^{iyx}\,dx
 &=\frac1{\sqrt{2\pi}}\int_0^\infty
 f(x)\bigl(e^{iyx}+e^{-iyx}\bigr)dx\\
 &=\frac2{\sqrt{2\pi}}\int_0^\infty f(x)\cos(xy)dx
 =\sqrt{\frac2\pi}\int_0^\infty f(x)\cos(xy)dx,
\end{align*}
which is \eqref{eq:even-transform-complete}.  If condition (B) holds with
$b>0$, then for $T>0$
\[
 |f(T)|=\sqrt{2\pi}|F(T)|
 =O\!\left(e^{-(1/2+b)T}\right).
\]
Therefore $T^{1/2}|f(T)|\to0$.  Guinand's Theorem~$3'$ explicitly lists this
limit as a sufficient condition for deleting the boundary term
\cite[pp.~112--115, Thms.~3 and~$3'$]{guinand1948}.
\end{proof}

\subsection{The rational trivial-character fibre}

\begin{theorem}[Termwise $k=\Q$, $\chi=1$ specialization]
\label{thm:weil-Q-specialization-complete}
In Weil's formula~(11), choose $k=\Q$ and the trivial character.  The retained
data are
\begin{equation}
 d=1,\quad\Delta=1,\quad N\mathfrak f=1,\quad A=(2\pi)^{-1},\quad
 (\eta_1,f_1,\varphi_1)=(1,0,0),\quad\delta_\chi=1.
 \label{eq:Q-fibre-data-complete}
\end{equation}
Let $F=\mathcal E(f)$ be in the intersection of
\cref{thm:guinand-weil-even-complete}.  Then Weil's finite-place term is
exactly
\begin{equation}
 -\sqrt{\frac2\pi}\sum_p\sum_{m\geq1}
 \frac{\log p}{p^{m/2}}f(m\log p),
 \label{eq:Q-prime-complete}
\end{equation}
his pole integral is
\begin{equation}
 \sqrt{\frac2\pi}\int_0^\infty
 f(x)(e^{x/2}+e^{-x/2})\,dx,
 \label{eq:Q-pole-complete}
\end{equation}
his conductor constant is
\begin{equation}
 -\frac{f(0)}{\sqrt{2\pi}}\log(2\pi),
 \label{eq:Q-conductor-complete}
\end{equation}
and the remaining real-place term is
\begin{equation}
 -\frac1{\sqrt{2\pi}}\operatorname{PF}\!\int_\R
 f(|x|)\frac{e^{|x|/2}}{|e^x-e^{-x}|}\,dx.
 \label{eq:Q-arch-complete}
\end{equation}
The prime series and pole integrals converge absolutely under condition (B).
\end{theorem}

\begin{proof}
For $\Q$, prime ideals are the ordinary primes, $N(p^m)=p^m$, and the
trivial character contributes $1$ in both finite-place orientations.  Since
$F$ is even,
\[
 F(m\log p)+F(-m\log p)
 =\frac{2}{\sqrt{2\pi}}f(m\log p)
 =\sqrt{\frac2\pi}f(m\log p).
\]
Inserting this into Weil's negative finite-place sum proves
\eqref{eq:Q-prime-complete}.  Splitting the pole integral at zero, changing
$x$ to $-x$ in the negative half, and using evenness gives twice the positive
half; $2/\sqrt{2\pi}=\sqrt{2/\pi}$, proving
\eqref{eq:Q-pole-complete}.  Since $\log A=-\log(2\pi)$ and
$F(0)=f(0)/\sqrt{2\pi}$, \eqref{eq:Q-conductor-complete} follows.
The real-place data in \eqref{eq:Q-fibre-data-complete} give Weil's kernel
$K_{1,0}(x)=e^{|x|/2}/|e^x-e^{-x}|$; linearity of the PF distribution and
$F=f(|\cdot|)/\sqrt{2\pi}$ give \eqref{eq:Q-arch-complete}.

If condition (B) holds with $b>0$, then at $m\log p>0$,
$F(\pm m\log p)=O(p^{-m(1/2+b)})$.  After the outside factor $p^{-m/2}$,
the absolute prime series is bounded by a constant multiple of
$\sum_{p,m}(\log p)p^{-m(1+b)}$, which converges.  The $e^{x/2}$ pole
integrand is $O(e^{-bx})$ and the $e^{-x/2}$ part decays faster.  These are
literal substitutions into Weil's primary formula
\cite[Collected Papers II, pp.~57--58, eq.~(11)]{weil1952}; other fields,
characters, conductors, and complex places remain outside this selected
fibre.
\end{proof}

\subsection{Zero pairing and the residual archimedean identity}

\begin{theorem}[Symmetric cutoff versus positive ordinates]
\label{thm:zero-pairing-complete}
Assume RH for $\zeta$, count zeros with multiplicity, and let
$F=\mathcal E(f)$, so that
$\Phi_F(1/2+iy)=g(y)$ is even.  Then, for every cutoff $T$ that does not pass
through an ordinate,
\begin{equation}
 \sum_{|\gamma|<T}\Phi_F(\tfrac12+i\gamma)
 =2\sum_{0<\gamma<T}g(\gamma).
 \label{eq:zero-pairing-complete}
\end{equation}
There is no omitted central term because $\zeta(1/2)<0$.
For a different self-dual $L$-function with central-zero multiplicity $m_0$,
the corresponding formula contains the additional term $m_0g(0)$.
\end{theorem}

\begin{proof}
The functional equation and conjugation symmetry pair every positive ordinate
$\gamma$ with $-\gamma$ with equal multiplicity.  Evenness gives
$g(-\gamma)=g(\gamma)$, so each noncentral pair contributes $2g(\gamma)$.
It remains only to check the central locus.

The alternating eta series at $s=1/2$ converges, and its even partial sums are
\[
 \sum_{k=1}^N\left((2k-1)^{-1/2}-(2k)^{-1/2}\right)>0.
\]
Their limit is therefore strictly positive.  The eta continuation identity
$\eta(s)=(1-2^{1-s})\zeta(s)$ gives
\[
 \zeta(1/2)=\frac{\eta(1/2)}{1-\sqrt2}<0.
\]
Thus $1/2$ is not a zero, proving \eqref{eq:zero-pairing-complete}.  Guinand's
positive-ordinate convention and Weil's symmetric cutoff are retained from
their respective primary formulas
\cite[pp.~107--115]{guinand1948}
\cite[Collected Papers II, pp.~57--58]{weil1952}.  The $m_0g(0)$ exception is
the unpaired central contribution and cannot be removed by evenness.
\end{proof}

\begin{theorem}[Guinand--Weil residual archimedean identity]
\label{thm:guinand-weil-residual-complete}
Assume the hypotheses of \cref{thm:guinand-weil-even-complete}, RH, and the
literal summation conventions of both primary formulas.  Write
\begin{align*}
 P_f&=\sum_p\sum_{m\geq1}\frac{\log p}{p^{m/2}}f(m\log p),\\
 A_+&=\int_0^\infty f(x)e^{x/2}dx,\qquad
 A_-=\int_0^\infty f(x)e^{-x/2}dx,\\
 H_f&=\frac12\int_0^\infty f(x)
 \left(\frac1x-\frac{e^{-x/2}}{\sinh x}\right)dx,\\
 Z_g(T)&=\sum_{0<\gamma<T}g(\gamma),\qquad
 L_g(T)=\int_0^Tg(t)\log\frac{t}{2\pi}dt,\\
 K(x)&=\frac{e^{|x|/2}}{|e^x-e^{-x}|}.
\end{align*}
Then both limits $Z_g=\lim_{T\to\infty}Z_g(T)$ and
$L_g=\lim_{T\to\infty}L_g(T)$ exist on this common subspace, and
\begin{equation}
 2A_--\operatorname{PF}\!\int_\R f(|x|)K(x)dx
 -f(0)\log(2\pi)
 =2H_f+\sqrt{\frac2\pi}L_g.
 \label{eq:guinand-weil-residual-complete}
\end{equation}
This is equality of scalar functionals on the stated common test subspace.  It
does not replace Weil's PF convention by an ordinary integral.
\end{theorem}

\begin{proof}
By \cref{thm:zero-pairing-complete}, Weil's symmetric zero limit is
$2\lim_T Z_g(T)$, so Weil's formula proves that $\lim_T Z_g(T)$ exists.
By \cref{thm:guinand-weil-even-complete}, Guinand's boundary term vanishes,
and his Theorem~$3'$ proves that
$\lim_T\{Z_g(T)-L_g(T)/(2\pi)\}$ exists.  Therefore
\[
 L_g(T)=2\pi\left[Z_g(T)-
 \left\{Z_g(T)-\frac{L_g(T)}{2\pi}\right\}\right]
\]
also has a limit.  Denote the two limits by $Z_g$ and $L_g$.
Guinand's theorem now states, in the notation above,
\begin{equation}
 P_f-A_+-H_f
 =-\sqrt{2\pi}\,Z_g+\frac1{\sqrt{2\pi}}L_g.
 \label{eq:guinand-scalars-complete}
\end{equation}
Indeed, his density term is $(2\pi)^{-1}L_g$, and
$\sqrt{2\pi}/(2\pi)=1/\sqrt{2\pi}$
\cite[pp.~114--115, Thms.~3 and~$3'$]{guinand1948}.

By \cref{thm:weil-Q-specialization-complete,thm:zero-pairing-complete}, Weil's
formula~(11) becomes
\begin{align}
 2Z_g={}&\sqrt{\frac2\pi}(A_++A_-)
 -\sqrt{\frac2\pi}P_f
 -\frac{f(0)}{\sqrt{2\pi}}\log(2\pi)\notag\\
 &-\frac1{\sqrt{2\pi}}
 \operatorname{PF}\!\int_\R f(|x|)K(x)dx.
 \label{eq:weil-scalars-complete}
\end{align}
Multiply \eqref{eq:weil-scalars-complete} by $\sqrt{2\pi}$:
\begin{equation}
 2\sqrt{2\pi}Z_g
 =2A_++2A_--2P_f-f(0)\log(2\pi)
 -\operatorname{PF}\!\int_\R f(|x|)K(x)dx.
 \label{eq:weil-scaled-complete}
\end{equation}
Solving \eqref{eq:guinand-scalars-complete} for $\sqrt{2\pi}Z_g$ and doubling
gives
\begin{equation}
 2\sqrt{2\pi}Z_g
 =-2P_f+2A_++2H_f+\frac{2}{\sqrt{2\pi}}L_g.
 \label{eq:guinand-scaled-complete}
\end{equation}
Equate \eqref{eq:weil-scaled-complete} and
\eqref{eq:guinand-scaled-complete}.  The terms $-2P_f+2A_+$ cancel exactly;
moving the rest gives \eqref{eq:guinand-weil-residual-complete}, since
$2/\sqrt{2\pi}=\sqrt{2/\pi}$.  Every cancellation is displayed, and the PF
term remains Weil's literal distribution
\cite[Collected Papers II, pp.~56--58]{weil1952}.  A direct kernel-level proof
independent of the two primary formulas would be stronger and remains a
separate obligation.
\end{proof}

\subsection{Montgomery's Poisson kernel as a Weil test}

\begin{theorem}[Exact Poisson-kernel test and proved term coordinates]
\label{thm:montgomery-weil-poisson-complete}
Let $1<\sigma<2$, $x\geq1$, $t\in\R$, and put
\[
 a=\sigma-\frac12,\qquad u=\log x,\qquad
 F_{a,t,u}(v)=e^{-a|v-u|}e^{-it(v-u)}.
\]
Then $F_{a,t,u}$ satisfies Weil's (A) and (B): it is continuous with one
permitted derivative jump at $v=u$, and for every
$0<b<a-1/2=\sigma-1$ it has the required exponential decay.  In the strip
$|\Re(s)-1/2|<a$,
\begin{equation}
 \Phi_{a,t,u}(s)
 =\frac{2a\,x^{s-1/2}}
 {a^2-(s-1/2-it)^2}.
 \label{eq:poisson-transform-complete}
\end{equation}
Consequently
\begin{equation}
 \Phi_{a,t,u}(\tfrac12+i\gamma)
 =\frac{(2\sigma-1)x^{i\gamma}}
 { (\sigma-\tfrac12)^2+(t-\gamma)^2},
 \label{eq:poisson-zero-complete}
\end{equation}
which is Montgomery's exact zero kernel.  At a positive integer $n$,
\begin{align}
 n^{-1/2}F(\log n)
 &=\begin{cases}
 x^{1/2-\sigma+it}n^{\sigma-1-it},&n\leq x,\\
 x^{\sigma-1/2+it}n^{-\sigma-it},&n>x,
 \end{cases}\label{eq:poisson-positive-prime-complete}\\
 n^{-1/2}F(-\log n)
 &=x^{1/2-\sigma+it}n^{-\sigma+it}.
 \label{eq:poisson-negative-prime-complete}
\end{align}
Thus the positive-log Weil prime term recovers Montgomery's two prime ranges,
while the negative-log term equals
\begin{equation}
 x^{1/2-\sigma+it}\frac{\zeta'}{\zeta}(\sigma-it)
 \label{eq:poisson-logderivative-complete}
\end{equation}
after Weil's outside minus sign.  Finally,
\begin{equation}
 \Phi_{a,t,u}(1)
 =\frac{x^{1/2}(2\sigma-1)}
 { (\sigma-1+it)(\sigma-it)},
 \label{eq:poisson-pole-complete}
\end{equation}
Montgomery's rational pole term.  The remaining gamma/trivial-zero
commutator is not claimed proved.
\end{theorem}

\begin{proof}
Away from $v=u$, differentiation multiplies $F$ by $a-it$ on the left and by
$-a-it$ on the right, so $F$ is continuous with one first-kind derivative
jump.  Assigning $F'(u)=-it$ gives the average of the two one-sided derivative
values, exactly as required by Weil's convention.  Its absolute value is
$e^{-a|v-u|}$; choosing
$0<b<a-1/2$ proves condition (B) for both $F$ and its one-sided derivatives.

Set $w=v-u$ and $z=s-1/2-it$.  For $|\Re z|<a$, absolute convergence permits
the split at $w=0$:
\begin{align*}
 \Phi(s)
 &=x^{s-1/2}\int_\R e^{-a|w|}e^{zw}dw\\
 &=x^{s-1/2}
 \left(\int_0^\infty e^{-(a-z)w}dw
       +\int_0^\infty e^{-(a+z)w}dw\right)\\
 &=x^{s-1/2}\left(\frac1{a-z}+\frac1{a+z}\right)
 =\frac{2a\,x^{s-1/2}}{a^2-z^2},
\end{align*}
proving \eqref{eq:poisson-transform-complete}.  At
$s=1/2+i\gamma$, $z=i(\gamma-t)$ and $2a=2\sigma-1$, giving
\eqref{eq:poisson-zero-complete}.

If $n\leq x$, then $\log n-u\leq0$, so direct substitution gives
\[
 F(\log n)=e^{-a(u-\log n)}e^{-it(\log n-u)}
 =x^{-a+it}n^{a-it}.
\]
Multiplication by $n^{-1/2}$ and $a=\sigma-1/2$ gives the first case of
\eqref{eq:poisson-positive-prime-complete}.  If $n>x$, the same calculation
with $|\log n-u|=\log n-u$ gives the second case.  Since $u\geq0$,
$-\log n-u\leq0$ for every $n\geq1$, giving
\eqref{eq:poisson-negative-prime-complete}.  For $\Re(\sigma-it)>1$,
\[
 -\frac{\zeta'}{\zeta}(\sigma-it)
 =\sum_{n\geq1}\Lambda(n)n^{-\sigma+it}.
\]
Weil's finite-place term has an outside minus sign, so its negative-log part
is \eqref{eq:poisson-logderivative-complete}.

At $s=1$, factor the denominator:
\begin{align*}
 a^2-(\tfrac12-it)^2
 &=(a-\tfrac12+it)(a+\tfrac12-it)\\
 &=(\sigma-1+it)(\sigma-it),
\end{align*}
which proves \eqref{eq:poisson-pole-complete}.  Comparing
\eqref{eq:poisson-zero-complete},
\eqref{eq:poisson-positive-prime-complete}, and
\eqref{eq:poisson-pole-complete} with Montgomery's complete formula gives
the advertised exact matches
\cite[pp.~185--187, eqs.~(18)--(22)]{montgomery1973}.  Still to be proved is
that $\Phi(0)$, the conductor constant, the PF real-place term, and
\eqref{eq:poisson-logderivative-complete} combine through the functional
equation into Montgomery's $-x^{1/2-\sigma+it}\zeta'/\zeta(1-\sigma+it)$ and
his entire trivial-zero series.  Until then, the full formulas are not
declared identical.
\end{proof}
